% O014 / D80 — Methods of Algebra, Volume 2 — id-ID
% Sumber: appendix2.tex, baris 502--605 (inklusif)
% ID stabil unit: o014.aljabr2.appendix-b.ind-abelian
\section{Ind-isasi kategori abelian}\label{sec:Indization-Abel}
Bagian ini melanjutkan pembahasan dan asumsi-asumsi terkait dari \S\S\ref{sec:ind-pro}---\ref{sec:Indization-functor}.

\begin{theorem}\label{prop:Ind-Abel}
	Jika $\mathcal{C}$ kategori abelian, maka $\IndC\mathcal{C}$ juga kategori abelian dan $\mathcal{C}\to\IndC\mathcal{C}$ merupakan funktor eksak setia-penuh.
	
	Secara dual, jika $\mathcal{C}$ kategori abelian, maka $\mathcal{C}\to\ProC\mathcal{C}$ juga merupakan funktor eksak setia-penuh di antara kategori abelian.
\end{theorem}
\begin{proof}
	Mula-mula tangani versi $\IndC\mathcal{C}$. Di bawah ini tuliskan $\Hom:=\Hom_{\IndC\mathcal{C}}$. Karena $\mathcal{C}\to\IndC\mathcal{C}$ mempertahankan $\varinjlim$ hingga (Lema \ref{prop:Ind-colim}) dan $\varprojlim$ hingga (Lema \ref{prop:Ind-lim}), funktor ini memetakan objek nol ke objek nol. Ini juga menunjukkan bahwa $\IndC\mathcal{C}$ mempunyai $\varinjlim$ dan $\varprojlim$ hingga.
	
	Selanjutnya, morfisme kanonik pada \eqref{eqn:coprod-prod-delta},
	\[ \delta: Y \sqcup Y \to Y \times Y, \quad Y \in \Obj(\IndC\mathcal{C}), \]
	selalu merupakan isomorfisme: berdasarkan konstruksi ``indeks demi indeks'' dari produk atau koproduk objek-ind, verifikasinya langsung tereduksi pada kasus $X,Y\in\Obj(\mathcal{C})$. Karena itu, mengikuti metode Proposisi \ref{prop:additive-cat-addition}, kita dapat mendefinisikan penjumlahan pada $\Hom(X,Y)$ dengan menetapkan $f+g$ sebagai komposisi
	\[ X \to X \times X \xrightarrow{f \times g} Y \times Y \xrightarrow{\delta^{-1}} Y \sqcup Y \to Y. \]
	Pembaca dapat memverifikasi bahwa, melalui
	\[ \Hom(X,Y)\simeq\varprojlim_i\varinjlim_j\Hom_{\mathcal{C}}(X_i,Y_j), \]
	operasi ini menjadi penjumlahan pada $\Hom_{\mathcal{C}}$ yang dicirikan dengan cara sama (ingat bahwa $\varinjlim_j$ terfilter). Jadi operasi tersebut benar-benar menjadikan $\IndC\mathcal{C}$ kategori-$\cate{Ab}$, kemudian kategori aditif, dan menjadikan $\mathcal{C}\to\IndC\mathcal{C}$ funktor aditif.
	
	Tinjau sebarang morfisme $f:X\to Y$ dalam $\IndC\mathcal{C}$. Kita menyatakan bahwa morfisme kanonik $\Coim(f)\to\Image(f)$ selalu merupakan isomorfisme. Menurut Lema \ref{prop:ind-object-morphism} (ii), kita boleh mengandaikan bahwa $f$ diinduksi oleh suatu keluarga morfisme $f_i:X_i\to Y_i$ dalam $\mathcal{C}$. Dalam kategori aditif, citra dan kocitra dideskripsikan melalui kernel dan kokernel (Proposisi \ref{prop:Im-Coim-additive}); ingat pula sifat-sifat yang diperoleh dalam pembuktian sebelumnya, seperti $\Ker(f)=\Yinjlim\Ker(f_i)$. Dengan ini, pernyataan langsung tereduksi pada $\Coim(f_i)\rightiso\Image(f_i)$. Maka $\IndC\mathcal{C}$ adalah kategori abelian.
	
	Terakhir, keeksakan $\mathcal{C}\to\IndC\mathcal{C}$ mengikuti karena funktor tersebut mempertahankan $\varinjlim$ dan $\varprojlim$ hingga.
	
	Konsep kategori abelian bersifat swadual (Proposisi \ref{prop:abelian-cat-duality}) dan $\ProC\mathcal{C}\simeq\IndC(\mathcal{C}^{\opp})^{\opp}$. Karena itu, pernyataan dual mengenai $\ProC\mathcal{C}$ langsung diperoleh.
\end{proof}

Perhatikan bahwa kategori abelian merupakan jenis khusus kategori aditif, dan keaditifan adalah sifat suatu kategori, bukan struktur tambahan; lihat pembahasan setelah Korolari \ref{prop:automatic-additivity}.

\begin{remark}\label{rem:Ind-k-linear}
	Jika lebih lanjut $\mathcal{C}$ diasumsikan sebagai kategori abelian $\Bbbk$-linear, dengan $\Bbbk$ gelanggang komutatif, maka $\IndC\mathcal{C}$ juga demikian dan $\mathcal{C}\to\IndC\mathcal{C}$ merupakan funktor $\Bbbk$-linear. Untuk melihatnya, pokoknya ialah mendefinisikan secara kanonik struktur modul-$\Bbbk$ pada setiap $\Hom(X,Y)$ sedemikian sehingga, pada ruas kanan bijeksi
	\[ \Hom(X, Y) \simeq \varprojlim_i \varinjlim_j \Hom_{\mathcal{C}}(X_i, Y_j), \]
	struktur itu tercermin sebagai struktur modul-$\Bbbk$ pada setiap $\Hom_{\mathcal{C}}(X_i,Y_j)$. Syarat ini dapat diambil sebagai definisi, tetapi kita harus menunjukkan bahwa ia hanya bergantung pada $X$ dan $Y$, bukan pada pilihan $(X_i)_i$ dan $(Y_j)_j$. Verifikasi ini tidak mengandung kesulitan esensial; rinciannya diserahkan sebagai latihan.
\end{remark}

\begin{lemma}
	Tinjau funktor $F:\mathcal{C}\to\mathcal{D}$ di antara kategori abelian. Funktor $\IndC F:\IndC\mathcal{C}\to\IndC\mathcal{D}$ dari Definisi--Proposisi \ref{def:Ind-functor} juga merupakan funktor aditif.
\end{lemma}
\begin{proof}
	Tinjau objek-ind $X=\Yinjlim X_i$ dan $Y=\Yinjlim Y_j$ di atas $\mathcal{C}$. Dari konstruksinya, pemetaan yang diinduksi $\IndC F$ pada tingkat morfisme sama dengan
	\[ \varprojlim_i \varinjlim_j \Hom_{\mathcal{C}}(X_i, Y_j) \xrightarrow{\text{diinduksi oleh $F$}} \varprojlim_i \varinjlim_j \Hom_{\mathcal{D}}(FX_i, FY_j). \]
	Dengan mengingat struktur penjumlahan pada himpunan-$\Hom$ dari $\IndC\mathcal{C}$ dan $\IndC\mathcal{D}$ (lihat pembuktian Teorema \ref{prop:Ind-Abel}), pemetaan di atas merupakan homomorfisme grup aditif.
\end{proof}

\begin{lemma}
	Jika kategori abelian $\mathcal{C}$ mempunyai $\varinjlim$ kecil terfilter, maka $\sigma:\IndC\mathcal{C}\to\mathcal{C}$ dari Proposisi \ref{prop:Ind-adjunction} merupakan funktor aditif.
\end{lemma}
\begin{proof}
	Hal ini dapat diverifikasi langsung seperti sebelumnya, atau diperoleh dari fakta berikut: $\sigma$ adalah adjoin kiri dari $\IndC\mathcal{C}\to\mathcal{C}$, sehingga Korolari \ref{prop:automatic-additivity} (v) menyiratkan keaditifannya.
\end{proof}

Jika $\mathcal{C}$ dan $\mathcal{D}$ bersifat $\Bbbk$-linear, dengan $\Bbbk$ gelanggang komutatif, dan $F$ juga funktor $\Bbbk$-linear, hasil-hasil di atas mempunyai generalisasi yang jelas.

Hasil berikut digunakan dalam \S\ref{sec:Tannaka-coend}.

\begin{proposition}\label{prop:Ind-ext-exactness}
	Tinjau funktor eksak kanan $F:\mathcal{C}\to\mathcal{D}$ di antara kategori abelian. Andaikan $\mathcal{D}$ mempunyai $\varinjlim$ kecil terfilter.
	\begin{enumerate}[(i)]
		\item Perluasan $\hat{F}:\IndC\mathcal{C}\to\mathcal{D}$ yang diberikan oleh Definisi--Proposisi \ref{def:Ind-hat-functor} merupakan funktor eksak kanan.
		\item Jika, lebih lanjut, $F$ merupakan funktor eksak dan $\varinjlim$ kecil terfilter bersifat eksak dalam $\mathcal{D}$, maka $\hat{F}$ merupakan funktor eksak.
		\item Dalam keadaan (i), jika $\mathcal{C}$ juga disyaratkan sebagai kategori kecil, maka $\hat{F}$ mempertahankan semua $\varinjlim$ kecil.
	\end{enumerate}
\end{proposition}
\begin{proof}
	Dua lema sebelumnya menunjukkan bahwa $\hat{F}$ adalah funktor aditif, sedangkan Definisi--Proposisi \ref{def:Ind-hat-functor} menunjukkan bahwa $\hat{F}$ mempertahankan $\varinjlim$ hingga. Jadi (i) terbukti.
	
	Sekarang andaikan hipotesis (ii) berlaku. Tujuan kita tereduksi pada pembuktian bahwa $\hat{F}$ mempertahankan $\Ker$. Misalkan $f:X\to Y$ morfisme dalam $\IndC\mathcal{C}$. Berdasarkan Lema \ref{prop:ind-object-morphism} (ii), kita boleh mengandaikan $X=\Yinjlim X_i$, $Y=\Yinjlim Y_i$, dan $f=\Yinjlim f_i$, dengan $f_i:X_i\to Y_i$ memenuhi syarat kompatibilitas. Terapkan lebih dahulu $F$, lalu $\varinjlim_i$, kepada barisan eksak
	\[ 0 \to \Ker(f_i) \to X_i \xrightarrow{f_i} Y_i. \]
	Hasilnya adalah barisan eksak
	\[ 0 \to \varinjlim_i F\Ker(f_i) \to \varinjlim_i FX_i \to \varinjlim_i FY_i. \]
	
	Menurut pembuktian Lema \ref{prop:Ind-lim}, $\Yinjlim\Ker(f_i)=\Yinjlim\Ker(f_i,0)$ memberikan $\Ker(f)=\Ker(f,0)$; karena itu suku pertama barisan tersebut teridentifikasi dengan $\hat{F}(\Ker(f))$. Di sisi lain, morfisme pada bagian kedua teridentifikasi dengan $\hat{F}f:\hat{F}X\to\hat{F}Y$. Ini membuktikan (ii).
	
	Sekarang tinjau (iii). Proposisi \ref{prop:Ind-hat-small} menyatakan bahwa $\hat{F}$ mempertahankan semua $\varinjlim$ kecil terfilter.\footnote{Catatan penerjemah: sumber menambahkan penunjuk ``(ii)'', padahal Proposisi~\ref{prop:Ind-hat-small} tidak terbagi menjadi butir-butir.} Bersama dengan keeksakan kanan, ini menyiratkan bahwa $\hat{F}$ mempertahankan semua $\varinjlim$ kecil.
\end{proof}

Kita lanjutkan dengan memusatkan perhatian pada ind-isasi kategori kecil. Kita memerlukan teori kategori Grothendieck dari \S\ref{sec:Grothendieck-cat}.

\begin{lemma}\label{prop:Ind-generator}
	Jika $\mathcal{C}$ kategori kecil yang mempunyai $\varinjlim$ hingga, maka $\IndC\mathcal{C}$ mempunyai pembangkit.
\end{lemma}
\begin{proof}
	Kita tahu bahwa $\IndC\mathcal{C}$ kokomplet (Lema \ref{prop:Ind-colim}) dan $\mathcal{C}$ kecil. Karena itu, dalam $\IndC\mathcal{C}$ kita dapat membentuk
	\[ s := \coprod_{X \in \Obj(\mathcal{C})} X. \]
	
	Kita tunjukkan bahwa $s$ adalah pembangkit. Untuk sepasang morfisme $f,g:X\rightrightarrows Y$ dalam $\IndC\mathcal{C}$, tuliskan $X=\Yinjlim X_i$. Berdasarkan konstruksi $s$, untuk setiap $i$ dapat dipilih $\epsilon_i$ sedemikian sehingga komposisi
	\[ X_i \xrightarrow{\text{morfisme kanonik koproduk}} s \xrightarrow{\epsilon_i} X \]
	sama dengan morfisme kanonik $X_i\to X$. Jika $f\epsilon=g\epsilon$ untuk setiap morfisme $\epsilon:s\to X$, maka dengan mengambil $\epsilon=\epsilon_i$ kita melihat bahwa tarikan balik $f$ dan $g$ ke setiap $X_i$ sama; jadi $f=g$.
\end{proof}

\begin{theorem}\label{prop:Ind-Grothendieck}
	Jika $\mathcal{C}$ kategori abelian kecil, maka $\IndC\mathcal{C}$ merupakan kategori Grothendieck.
\end{theorem}
\begin{proof}
	Teorema \ref{prop:Ind-Abel} menyatakan bahwa $\IndC\mathcal{C}$ kategori abelian. Kita verifikasi satu demi satu syarat kategori Grothendieck.
	\begin{description}
		\item[Kokomplet] Termuat dalam Lema \ref{prop:Ind-colim}.
		\item[Pembangkit] Keberadaannya dijamin oleh Lema \ref{prop:Ind-generator}.
		\item[Keeksakan $\varinjlim$ kecil terfilter] Misalkan $I$ kategori kecil terfilter, $\alpha,\beta,\gamma:I\to\IndC\mathcal{C}$ funktor, dan $\alpha\to\beta\to\gamma$ morfisme sedemikian sehingga
		\[ 0 \to \alpha(i) \to \beta(i) \to \gamma(i) \to 0 \]
		merupakan barisan eksak pendek untuk setiap $i\in\Obj(I)$, kita hendak membuktikan bahwa
		\[ 0 \to \varinjlim \alpha \to \varinjlim \beta \to \varinjlim \gamma \to 0 \]
		juga eksak. Seperti dijelaskan dalam Contoh \ref{eg:limit-exactness}, $\varinjlim$ selalu mempertahankan $\Coker$; pokoknya adalah membuktikan bahwa ia mempertahankan $\Ker$. Namun, $\Ker$ merupakan suatu $\varprojlim$ hingga, sehingga persoalan ini diselesaikan oleh bagian kedua Lema \ref{prop:Ind-lim}.
	\end{description}
	Ini membuktikan pernyataan.
\end{proof}
